\documentclass{article}
\usepackage{/globals/anki}
\ankiprefix{daten}
\ankitopic{Daten}
\ankishow
\begin{document} 
Daten sind eine Menge an Zeichen in einem bestimmtem Syntax. Daten mit Kontext sind Informationen. Informationen mit Semantik ist Wissen.
\ankinote{daten}{Daten}{Daten = Zeichen + Syntax}
\ankinote{information}{Informationen}{Informationen = Daten + Kontext}
\ankinote{wissen}{Wissen}{Wissen = Information + Semantik}

Daten können in mehrere Verwendungszwecke kategorisiert werden.
\begin{description}
 \item[Stammdaten] Daten die sich nicht häufig ändern.
 \item[Bewegungsdaten] Daten die sich häufiger ändern.
 \item[Transferdaten] Daten die gerade zwischen Systemen übertragen werden.
 \item[Vormerkdaten] Temporäre Daten.
\end{description}
\ankinote{arten}{Arten}{Stamm, Bewegungs, Transfer, Vormerk}
\ankinote{stamm}{Stammdaten}{Seltene Änderung}
\ankinote{bewegungs}{Bewegungsdaten}{Häufige Änderung}
\ankinote{transfer}{Transferdaten}{Übertragung}
\ankinote{vormerk}{Vormerkdaten}{Temporär}
\end{document}